% pdflatex myfile.tex
% makeindex myfile.idx
% pdflatex myfile.tex
% biber myfile
% pdflatex myfile.tex
% pdflatex myfile.tex

\documentclass[
    12pt,
    ngerman,
    a4paper,
    oneside
]{book}

% UTF-8 + T1 Fontencoding
\usepackage[utf8]{inputenc}
\usepackage[T1]{fontenc}
\usepackage{helvet}

% Import von Meta-Daten (Autor, Titel, Prüfer, etc.)
\newcommand{\titel}{Selbstständige Inbetriebnahme und Verwaltung einer Webanwendung zur Verwaltung pädagogischer Materialien durch Jugendliche}
\newcommand{\untertitel}{}
\newcommand{\untertitelDeckblatt}{}
\newcommand{\art}{Exposé}
\newcommand{\fachgebiet}{zur Erlangung des akademischen Grades\\ Bachelor of Science (B.\,Sc.) im Studienfach\xspace}
\newcommand{\autor}{Torben Peter Wager}
\newcommand{\keywords}{Bachelorarbeit, Torben Peter Wager}
\newcommand{\studienbereich}{Informatik\xspace}
\newcommand{\matrikelnr}{937525}
\newcommand{\erstgutachter}{Corina Kopka}
\newcommand{\zweitgutachter}{Eike Petersen}
\newcommand{\jahr}{2026}
\newcommand{\hochschule}{HAW Kiel}
\newcommand{\ort}{Kiel}
\newcommand{\logo}{HAWLogo.pdf}
\newcommand{\creator}{Tex Live 2026}

% Import von ausgelagerten Packages
% Hier könnnen während der Bearbeitung der Arbeit weitere Packages eingebunden werden.

\usepackage{nomencl}
\usepackage{makeidx}

% sorgt dafür das generierter Text auf deutsch ist (Bsp. Inhaltsverzeichnis)
\usepackage[ngerman]{babel}

% für Tabellen
\usepackage{array}

% zusätzliche Mathesymbole
\usepackage{amsmath,amssymb}
\usepackage{bm}

% Calibri esque Font
\usepackage[sfdefault]{carlito}

% FÜr die Aufzählungen
\usepackage{enumitem}

% zum generieren von Literaturverzeichnis
\usepackage{csquotes}
\usepackage{biblatex}
% Robuste Zitier-Wrapper für bewegte Argumente (Captions, etc.)
\DeclareRobustCommand{\citerobust}[1]{\cite{#1}}

% Kapitel Formatierung anpassen
\usepackage{titlesec}
\titleformat{\chapter}[hang]
  {\normalfont\huge\bfseries}
  {\thechapter}               
  {1em}                       
  {}    

% Standard Kopfzeile
\usepackage{fancyhdr}
\fancypagestyle{fancy}{
  \fancyhf{}
  \rhead{\rightmark}
  \cfoot{\thepage}
  \renewcommand{\headrulewidth}{0.4pt}
}

% für keine Header
\fancypagestyle{plain}{
  \fancyhf{}
  \cfoot{\thepage}
  \renewcommand{\headrulewidth}{0pt}
}

% sorgt dafür, dass Leerzeichen hinter parameterlosen Makros nicht als Makroendezeichen interpretiert werden
\usepackage{xspace}

% zum Laden von Grafiken
\usepackage{graphicx}

%Zum färben von Text
\usepackage{xcolor}

% zum setzen von Ankern
\usepackage{hyperref}

% Standardisiertes Abbildungsformat
\usepackage[font=small,labelfont=bf,labelsep=colon]{caption}
\newcommand{\standardfigure}[4][]{%
  \begin{figure}[htbp]
    \centering
    \includegraphics[width=0.85\textwidth]{#2}
    \caption[#3]{#4}
    \label{fig:#1}
  \end{figure}
}

% Laden der .bib
\addbibresource{Bibliography.bib}

\makenomenclature{}

\renewcommand{\nomname}{Abkürzungen und Glossar}
\setlength{\nomlabelwidth}{3cm}         % Platz für das Kürzel
\setlength{\nomitemsep}{0.5\baselineskip} % Abstand zwischen den Einträgen
\renewcommand{\nomlabel}[1]{\textbf{#1}\dotfill}

\begin{document}

\thispagestyle{plain}

\begin{titlepage}
\vspace*{\fill}
\begin{flushleft}


\large{\textbf{\titel}}\\[7.5ex]
\Large{\textbf{\art}}\\[1.5ex]

\large{\fachgebiet \studienbereich}\\[12ex]

%\includegraphics[width=0.7\textwidth]{HAWLogo.pdf}\\[12ex]
\includegraphics[width=0.7\textwidth]{Abbildung/HAW_Kiel_Logo.pdf}\\[12ex]

\normalsize{}
\begin{tabular}{ll}
vorgelegt von:  & \quad \autor\\[1.2ex]
Matrikelnummer: & \quad \matrikelnr\\[1ex]
Erstgutachter:  & \quad \erstgutachter\\[1ex]
Zweitgutachter: & \quad \zweitgutachter\\[1ex]
eingereicht in: & \quad \ort, am xx.\,xxxxxx\,20xx
\end{tabular}

\end{flushleft}
\vspace*{\fill}
\end{titlepage}

\section*{Problemstellung}
Im Rahmen von Jugendfreizeiten oder Jugendtreffs werden häufig pädagogische Materialien wie Spielanleitungen, Theaterskripte oder Liedblätter erstellt und eingesetzt.
Diese Materialien werden meist wiederholt neu erstellt, da sie nicht gespeichert, oder gesammelt zugänglich gemacht werden.
Als Lösung für dieses Problem sollen die Materialien in einer Webanwendung gespeichert und verwaltet werden.
Diese Arbeit soll untersuchen, ob es ehrenamtlichen Mitarbeitenden zwischen 14 und 21 Jahren möglich ist, eine Webanwendung, welche ihnen bei der Organisation solcher Materialien hilft, selbstständig einzurichten und zu verwalten. 


\section*{Forschungsfrage}

\begin{quote}
Wie unterscheiden sich manuelle und containerisierte Installation und Wartung einer Webanwendung durch ehrenamtliche Jugendliche hinsichtlich Aufgabenerfolg, Bearbeitungszeit und benötigter Hilfestellung?
\end{quote}

\noindent Im Rahmen der Arbeit sollen zwei Bereitstellungsansätze für denselben Webapp-Prototypen verglichen werden: eine manuelle Installation von Webserver und Datenbank und eine containerisierte Bereitstellung mit Docker.
In beiden Ansätzen erhalten die Jugendlichen eine detaillierte Schritt-für-Schritt Anleitung und führen dieselben Aufgaben zur Inbetriebnahme und Aktualisierung auf einem Standardieserten Linux-System durch.
Die Qualitätsmerkmale nach denen die Evaluation quantifiziert werden soll:
\begin{enumerate}[
    label={-},
    leftmargin=*,
    itemsep=0.4em,
    topsep=0.4em
]
    \item \textbf{Aufgabenerfolg}: Anzahl der Aufgaben, die nach Abschluss des Versuchs und ohne Eingreifen richtig abgeschlossen werden.
    \item \textbf{Zeitaufwand}: Zeit vom Start der Bearbeitung einer Aufgabe bis zu ihrem Abschluss.
    \item \textbf{Benötigte Hilfestellung}: Frequenz und Dauer mit der Hilfestellung zum Lösen der Aufgaben in Anspruch genommen werden.
\end{enumerate}

\section*{Forschungsstand}
Vorhandene Forschung im Bereich digitaler Kompetenz Jugendlicher befasst sich überwiegend mit Nutzung von Suchmaschinen zur Informationsgewinnung oder sozialen Medien im Rahmen zwischenmenschlicher Interaktion.\cite{schaumburg2019nutzung}\cite{hartmann2020digitale}
Dienste wie die Jugendapp von jugendarbeit.digital oder GroupSenz sind Plattformen, welche Funktionen zur Verwaltung von Jugendangeboten bereitstellen, allerdings nicht auf die selbstständige Inbetriebnahme und Wartung durch Jugendliche ausgelegt sind.\cite{jugendarbeit.digital}\cite{groupsenz}
In den Recherchierten Arbeiten fanden sich keine Untersuchungen, die sich mit der Inbetriebnahme und Wartung einer Webanwendung durch Jugendliche befassten.
Auch die Verwendung von Containerisierung zur Einrichtung von digitalen Diensten durch Jugendliche ist kaum erforscht.
Es ist daher unklar, ob die Verwendung von Containerisierung im Vergleich zu einer manuellen Installation von Webserver und Datenbank die Inbetriebnahme und Wartung durch Jugendliche erleichtert.
Skriptbasierte Konfiguartion, virtuelle Maschinen, sowie Product- und Function-as-a-Service werden als alternative Ansätze eingeordnet, sind jedoch nicht Gegenstand dieser Arbeit, da sie entweder zu komplex oder zu simpel sind, um die Forschungsfrage zu beantworten.

\section*{Methodik}
Zu Beginn der Arbeit wird der Webanwendungsprototyp in einem agilen Entwicklungsprozess entwickelt.\cite{tremp2022agile}
In der Anforderungsanalyse werden zuerst die Stakeholdergruppen innerhalb der ehrenamtlichen Jugendarbeit identifiziert.
Zu unterscheiden sind die diversen Altersgruppen und verschiedenen Rollen innerhalb der Jugendarbeit.
In Interviews mit diesen Stakeholdergruppen werden danach Anforderungen an die Webanwendung definiert.
Diese werden als User Stories formuliert, die als Grundlage für die Entwicklung des Prototyps dienen.
Bei Durchführung der Usability-Tests sollen Testverfahren an die Arbeit mit Jugendlichen angepasst werden.\cite{Rose2018}
Damit die Entwicklungszeit im Hinblick auf Evaluation der Bereitstellungsansätze nicht zu lang wird, soll der Umfang der Webanwendung auf die Kernfunktionen zur Verwaltung von pädagogischen Materialien beschränkt werden.
Anhand dieser Grundanforderungen wurde sich für Django entschieden.
Das Framework stellt viele Funktionen für die Verwaltung von Datenbanken bereits bereit und ermöglicht damit eine schnellere Entwicklung des Prototypen.\cite{djangoDocs}

Nach abgeschlossener Entwicklung des Prototyps wird die Webanwendung in die Bereitstellungsansätze integriert.
Webserver und Datenbank werden in einem Docker-Container als Image bereitgestellt, während die manuelle Installation über git und Paketmanager erfolgt.
Für beide Bereitstellungsansätze werden Schritt-für-Schritt Anleitungen erstellt, die die Jugendlichen durch die Inbetriebnahme und Wartung der Webanwendung führen.
Mit fertigen Anleitungen werden Testpläne für die beiden Bereitstellungsansätze erstellt.
Eine Gruppe aus zwölf Ehrenamtlichen wird in vier Untergruppen unterteilt, welche die verschiedenen Bereitstellungsansätze in unterschiedlichen Reihenfolgen testen.
Eine Gruppe testet zuerst die manuelle Installation und danach die containerisierte Bereitstellung, eine Zweite tut dies umgekehrt.
Die beiden verbleibenden Gruppen testen jeweils nur eine der beiden Bereitstellungsansätze.
Die Aufgaben innerhalb der Testgruppen sind identisch und decken die gesammte Inbetriebnahme, sowie ausgewählte Wartungsaufgaben ab.
Die Ergebnisse werden anhand der Merkmale Aufgabenerfolg, Zeitaufwand und Benötigte Hilfestellung ausgewertet.
Diese Merkmale sollen quantifizierbar sein, um die beiden Bereitstellungsansätze miteinander vergleichen zu können.

Die Testleitung greift nur im Notfall ein, um bei Aufgaben zu helfen.
Als Aufgabenerfolg zählt eine Aufgabe, die ohne Eingreifen der Testleitung korrekt abgeschlossen wurde.
Der Zeitaufwand wird von Beginn der Bearbeitung einer Aufgabe bis zu ihrem Abschluss gemessen.
Die benötigte Hilfestellung wird von der Testleitung dokumentiert, indem sie festhält, wie oft und lange in der Anleitung nachgelesen werden musste, um die Aufgabe zu lösen.

\section*{Gliederung}
\begin{enumerate}[
    label=\arabic*.,
    leftmargin=*,
    itemsep=0.4em,
    topsep=0.2em
]
    \item \textbf{Einleitung} (10\%) \\
    Problemstellung, Forschungsfrage und Zielsetzung

    \item \textbf{Grundlagen und Forschungsstand} (15\%)\\
    Digitale Kompetenzen Jugendlicher und Bereitstellungsansätze für Webanwendungen

    \item \textbf{Methodik} (15\%)\\
    Beschreibung der Vorgehensweise, Probanden, Aufgaben, Testaufbau, Testdurchführung, Erhebungsmethoden und Auswerungsverfahren.
    
    \item \textbf{Anforderungsanalyse und Prototypentwicklung} (10\%)\\
    Erheben der Anforderungen an die Webanwendung und Entwicklung eines Prototyps.
    
    \item \textbf{Entwicklung der Bereitstellungsansätze} (10\%)\\
    Umsetzung der containerisierten und manuellen Bereitstellung sowie Erstellung der Anleitungen für die Testgruppen.
    
    \item \textbf{Entwurf und Durchführung der Tests} (20\%)\\
    Erstellen der Testpläne und Dokumentation der Ergebnisse und Beobachtungen während der Tests.
    
    \item \textbf{Diskussion} (15\%)\\
    Einordnen der Testergebnisse, Vergleich der Bereitstellungsansätze und Beantwortung der Forschungsfrage.
    
    \item \textbf{Fazit und Ausblick} (5\%)\\
    Zusammenfassung der Ergebnisse und Ausblick auf mögliche Weiterentwicklungen.


\end{enumerate}

\section*{Zeitplan}
Die Arbeitszeit umfasst 13 Wochen und soll sich grob in drei Phasen gliedern:
\begin{enumerate}[
    label=\Roman*.,
    leftmargin=*,
    itemsep=0.4em,
    topsep=0.4em
]
    \item Anforderungsanalyse und Prototypentwicklung (ca. 4 Wochen)
    \begin{enumerate}[label={-}, itemsep=0.1\baselineskip]
        \item Anforderungen der Webanwendung definieren (ca. 1 Woche)
        \item Prototyp entwickeln (ca. 3 Wochen)
    \end{enumerate}
    \item Erstellung der Testpläne und Durchführung der Tests (ca. 5 Wochen)
    \begin{enumerate}[label={-}, itemsep=0.1\baselineskip]
        \item Bereitstellungsansätze implementieren (ca. 1 Woche)
        \item User Guides und Testpläne für die Bereitstellungsansätze erstellen (ca. 2 Wochen)
        \item Usability-Tests durchführen (ca. 2 Wochen)
    \end{enumerate}
    \item Auswertung der Ergebnisse und Erstellung der Arbeit (ca. 4 Wochen)
    \begin{enumerate}[label={-}, itemsep=0.1\baselineskip]
        \item Ergebnisse auswerten (ca. 2 Wochen)
        \item Arbeit finalisieren (ca. 2 Wochen)
    \end{enumerate}
\end{enumerate}

\noindent Mit dem Schreiben soll bereits während der Prototypentwicklung begonnen und während der Tests laufend weitergeführt werden, um die Arbeit innerhalb des Zeitplans fertigzustellen.
Der Umfang der Thesis soll ohne Anhang circa 10.000 bis 15.000 Wörter betragen.

\section*{Literatur}
Für die methodische Arbeit während der Entwicklung der Webapp wird Literatur zu agiler Prototypentwicklung herangezogen.\cite{tremp2022agile}
Dort finden sich Hinweise zu Anforderungsanalyse, Usability-Tests und Evaluation von Softwareprototypen.
Zum Verfassen der Anleitungen ist Literatur zu didaktischen Methoden in der Arbeit mit Technik notwendig.\cite{wallwork2014user}
Durch diese wird sichergestellt, dass die Anweisungen für die Jugendlichen verständlich und anwendbar sind.
Die Entwicklung der Webanwendung, sowie die Bereitstellung dieser, wird auf Grundlage von Dokumentation der verwendeten Frameworks und Bibliotheken erfolgen.\cite{dockerDocs}\cite{djangoDocs}
Um die Arbeit in den Forschungsstand einzuordnen, werden Publikationen zu digitalen Kompetenzen Jugendlicher\cite{schaumburg2019nutzung}\cite{hartmann2020digitale} und zu bestehenden Plattformen für die Jugendarbeit\cite{jugendarbeit.digital}\cite{groupsenz} herangezogen.

% Literaturverzeichnis
\printbibliography{}

\end{document}